\uuid{iAKk}
\titre{Application linéaire}
\auteur{hanani ; mzari}
\organisation{exo7}
\licence{CC BY-NC-SA 4.0}
\datecreate{2020-01-20}
\niveau{L1}
\module{Algèbre}
\chapitre{Application linéaire}
\sousChapitre{Définition}
\difficulte{3}
\theme{Applications linéaires, Application linéaire}
\exo7qcmid{lille-2.applications-lineaires.007}

\contenu{
\begin{question}
\qclassification{107.01}
\qtags{motcle=Application linéaire}
On considère les deux applications suivantes :
$$\begin{array}{rccc}f:&\Rr^2&\to&\Rr\\
& (x,y)&\to & |x+y| \end{array} \quad \mbox{et} \quad \begin{array}{rccc}g:&\Rr^2&\to&\Rr^2\\
& (x,y)&\to &\big(\max (x,y)\, , \, \min (x,y)\, \big). \end{array}$$
Quelles sont les assertions vraies ?
\begin{answers}
\good{$f(1,-1)=0$}.
\bad{$f$ est une application linéaire}.
\good{$g(0,0)=(0,0)$}.
\bad{$g$ est une application linéaire}.
\end{answers}
\begin{explanations} $f$ n'est pas linéaire. Contre-exemple : $f((1,0)+(-1,0))=f(0,0)=0$ et $f(1,0)+f(-1,0)=|1|+|-1|=2$.
\vskip0mm
$g$ n'est pas linéaire. Contre-exemple :  $g((1,0)+(-1,0))=g(0,0)=(0,0)$
et $g(1,0)+g(-1,0)=(1,0)+(0,-1)=(1,-1)$.
\end{explanations}
\end{question}
}
